% !PRINTED 313
\weberpage{313}
%<<CONT
\weq{14}{x = y(\lambda+\mu z) = \frac{(\lambda+\mu z)\gamma_2}{z^2+3}}%
\footnote{Vgl. \gesperrt{Kiepert}, Auflösung der Gleichung 5ten Grades. Crelles Journal, Bd.\,87, S.\,114.}

sind, wenn $\lambda, \mu$ beliebige Parameter bedeuten, die für alle fünf
Werte $x$ die gleichen sind.

Setzt man diese Gleichung in die Form

\weq{15}{x^5 - 5ax^3 - 5b\gamma_2 x - bc\gamma_2^2 = 0,}

so ergeben sich nach einigen Rechnungen mit Benutzung der besonderen
Werte von $x$ für $\gamma_3=0$ und $\gamma_3=\infty$ für $a, b, c$ die
folgenden Ausdrücke:

\weq{16}{\begin{aligned}
a &= 8\lambda^3 - 72\lambda\mu^2 + \gamma_3(\lambda^2\mu - \mu^3),\\
b &= \lambda^4 + 18\lambda^2\mu^2 - 27\mu^4 + \gamma_3\lambda\mu^3,\\
c &= \lambda^5 + 10\lambda^3\mu^2 + 45\lambda\mu^3 + \gamma_3\mu^5. % [sic] Gedruckt ist μ^3.
\end{aligned}}

\par\noindent\rule{0.3\linewidth}{0.4pt}\par

Es soll noch eine Resolvente 5ten Grades mit Benutzung der Funktion
$f(\omega)$ gebildet werden.

Wir setzen:

\weq{17}{\begin{gathered}
u = f(\omega), \quad v_\infty = f(5\omega), \quad v_z = f\left(\frac{\omega+c}{5}\right)\\
c \equiv 0 \pmod{48}, \quad z \equiv c \pmod{5},
\end{gathered}}

und haben nach §\,73 zwischen $u$ und $v$ die Gleichung 6ten Grades:

\weq{18}{u^6 + v^6 - u^5v^5 + 4uv = 0.}

Wir untersuchen den Einfluß, den die drei Vertauschungen

\[
(\omega, \omega+2), \quad \left(\omega, -\frac{1}{\omega}\right), \quad
\left(\omega, \frac{\omega-1}{\omega+1}\right)
\]

auf die Größen $v_z$ haben. Durch diese Vertauschungen geht $c$ über in
$c_1, c_2, c'$, die nach §\,69, (9), (11) und §\,73, (6) durch die
Kongruenzen

\weq{19}{c_1 \equiv c+2, \quad c_2 \equiv -\frac{1}{c}, \quad c' \equiv \frac{c-1}{c+1} \pmod{5}}

bestimmt sind. Abgesehen von dieser Änderung des Index gehen nach §\,73,
(4) und (9) die $v$ über in

\weq{20}{e^{-\frac{5\pi i}{12}}v, \quad v, \quad -\frac{\sqrt{2}}{v}.}
